%% ============================================================
%% AULA 3 — Aprendendo a rede: Parte 1
%% GBC073 — Inteligência Computacional (tema InteligenciaComputacional)
%%
%% Adaptada da Lecture 3 do 11-785 (Intro to Deep Learning, CMU, S26):
%% "Learning the network: Part 1".
%% Ênfase moderna: diferenciabilidade, autograd, ERM e perdas em PyTorch.
%% Perceptron de Rosenblatt e ADALINE/MADALINE como perspectiva histórica.
%%
%% Compilar:  latexmk -xelatex aula03.tex   (XeLaTeX, duas passadas)
%% ============================================================
\documentclass[aspectratio=169, 11pt]{beamer}

\makeatletter
\def\input@path{{./}{../}}
\makeatother

\usetheme{InteligenciaComputacional}   % ANTES do babel: fontspec vem primeiro
\usepackage{babel}
\babelprovide[import, main]{brazilian}

% ---------- código Python/PyTorch --------------------------------
\usepackage{listings}
\lstdefinestyle{py}{
  language=Python,
  basicstyle=\ttfamily\footnotesize,
  keywordstyle=\color{icAzul}\bfseries,
  commentstyle=\color{icCinza}\itshape,
  stringstyle=\color{icVerde},
  emph={torch, nn, optim, backward, grad},
  emphstyle=\color{icMagenta},
  showstringspaces=false,
  columns=fullflexible,
  keepspaces=true,
  xleftmargin=1em,
}
\lstset{style=py}

% ---------- macro local desta aula -------------------------------
\newcommand{\dive}{\operatorname{div}}   % divergência rede x alvo

\title[Aula 3 — Aprendendo a rede I]{Aula 3: Aprendendo a rede — Parte 1}
\subtitle[GBC073 — Inteligência Computacional]{Do perceptron à minimização de risco empírico \textbullet\ por que diferenciabilidade é tudo}
\author[Prof.\ Marcelo Albertini]{Prof.\ Marcelo Keese Albertini}
\institute{GBC073 — Inteligência Computacional \textbullet\ Universidade Federal de Uberlândia}
\date{2026}

\begin{document}

% ------------------------------------------------------------ capa
\begin{frame}[plain, noframenumbering]
  \titlepage
\end{frame}

% ------------------------------------------------------------ roteiro
\begin{frame}{Roteiro da aula}
  \framesubtitle{Baseada na Lecture 3 do 11-785 (CMU) — com olhar de deep learning moderno}
  \begin{itemize}\setlength{\itemsep}{0.5ex}
    \item<1-> \textbf{O problema de aprendizado:} a rede é uma função paramétrica
          $f(\vx; W)$ — aprender $=$ estimar $W$ a partir de \emph{amostras}
    \item<2-> \textbf{Histórico:} a regra do perceptron — e por que ela
          \alert{não} treina redes de múltiplas camadas
    \item<3-> \textbf{A virada moderna:} ativações e perdas \alert{diferenciáveis},
          a decisão que tornou o deep learning possível
    \item<4-> \textbf{ERM:} o arcabouço que todo treinamento em PyTorch
          implementa até hoje
    \item<5-> \textbf{Ponte para a Aula 4:} treinar $=$ otimizar uma função
          diferenciável $\Rightarrow$ descida de gradiente
  \end{itemize}

  \onslide<5->{%
  \begin{ideiabox}
    Uma ideia atravessa a aula: só ajustamos bilhões de parâmetros porque cada
    pedaço da rede responde ``se eu mexer um pouquinho aqui, o erro muda quanto?''.
  \end{ideiabox}}
\end{frame}

% ============================================================
\section{O problema de aprendizado}

% ------------------------------------------------------------
\begin{frame}{Recapitulação: redes representam qualquer função}
  \begin{itemize}\setlength{\itemsep}{0.4ex}
    \item Redes neurais são \textbf{aproximadores universais}: funções booleanas,
          fronteiras de classificação e funções contínuas em compactos
    \item Desde que a arquitetura tenha \alert{capacidade suficiente}
          (neurônios/camadas em número adequado)
    \item Voz $\to$ transcrição, imagem $\to$ legenda, estado do jogo $\to$ jogada:
          todas são funções — logo, todas são modeláveis por uma rede
  \end{itemize}

  \begin{alertabox}
    ``Existe uma rede que representa $g$'' \alert{não} diz \emph{como encontrá-la}:
    aproximação universal é teorema de existência; aprendizado é \emph{busca}.
  \end{alertabox}

  \begin{ideiabox}
    Separar \emph{representação} de \emph{aprendizado}: a Aula 2 respondeu
    ``o que a rede \emph{pode} computar''; hoje, ``como fazê-la computar o que
    queremos''.
  \end{ideiabox}
\end{frame}

% ------------------------------------------------------------
\begin{frame}{A rede é uma função paramétrica}
  \begin{columns}[T]
    \begin{column}{0.52\textwidth}
      \begin{defbox}{Aprender a rede}
        Fixada a arquitetura, a rede computa $f(\vx;\, W)$, onde $W$ reúne
        \textbf{todos os pesos e vieses}. \emph{Aprender} é determinar $W$
        tal que $f(\cdot\,; W)$ aproxime a função-alvo $g$.
      \end{defbox}
      \begin{itemize}\setlength{\itemsep}{0.3ex}
        \item A arquitetura (quantas camadas, quem liga em quem) é
              \emph{decisão de projeto} — assumimos hoje que ela basta
        \item Os \textbf{parâmetros} $W$ são o que o treinamento ajusta
        \item Rede direta (\emph{feedforward}): sem laços — redes com
              realimentação ficam para o futuro
      \end{itemize}
    \end{column}
    \begin{column}{0.44\textwidth}
      \begin{center}
        \scalebox{0.62}{\mlp{3,4,4,2}}
      \end{center}
      \begin{center}
        {\small\color{icCinza} Cada seta carrega um peso $\pesoij{j}{i}$;
        o conjunto de todos eles é $W$.}
      \end{center}
    \end{column}
  \end{columns}
\end{frame}

% ------------------------------------------------------------
\begin{frame}{Construir a rede à mão não escala}
  \begin{columns}[T]
    \begin{column}{0.55\textwidth}
      \begin{itemize}\setlength{\itemsep}{0.4ex}
        \item Fronteira simples (losango): dá para \emph{projetar} os pesos —
              uma reta por lado, um E lógico no topo
        \item Reconhecer um dígito $28{\times}28$? Transcrever fala?
              \alert{Impossível} escrever os pesos à mão
        \item Precisamos de \textbf{estimação automática}: calcular $W$ que
              minimize a discordância entre $f(\vx; W)$ e $g(\vx)$
      \end{itemize}
      \begin{ideiabox}
        Ideal: minimizar o desacordo $f \times g$ em \emph{todo} o domínio.
        Mas para isso precisaríamos conhecer $g$ em todo lugar — e aí não
        haveria o que aprender.
      \end{ideiabox}
    \end{column}
    \begin{column}{0.41\textwidth}
      \begin{center}
      \scalebox{0.68}{%
      \begin{tikzpicture}
        \draw[->, icCinza] (-1.7,0) -- (1.7,0) node[below right] {\small $x_1$};
        \draw[->, icCinza] (0,-1.7) -- (0,1.7) node[above left] {\small $x_2$};
        \draw[thick, icMagenta, fill=icMagenta!12]
          (1.2,0) -- (0,1.2) -- (-1.2,0) -- (0,-1.2) -- cycle;
        \node[icMagenta] at (0,0.35) {\small $+1$};
        \node[icCinza] at (1.15,1.15) {\small $0$};
      \end{tikzpicture}}
      \end{center}
      \begin{center}
        {\small\color{icCinza} 4 retas $+$ 1 conjunção: fácil.
        Agora imagine $10^7$ pesos.}
      \end{center}
    \end{column}
  \end{columns}
\end{frame}

% ------------------------------------------------------------
\begin{frame}{Amostrar a função: os dados de treinamento}
  \begin{defbox}{Conjunto de treinamento}
    Não conhecemos $g$ inteira; conhecemos \textbf{amostras} dela:
    \[
      \D = \{(\vx_1, d_1),\, (\vx_2, d_2),\, \ldots,\, (\vx_N, d_N)\},
      \qquad d_i = g(\vx_i).
    \]
    Ex.: imagens e seus rótulos; áudios e suas transcrições.
  \end{defbox}

  \begin{itemize}\setlength{\itemsep}{0.3ex}
    \item Ajustamos $W$ para que $f(\vx_i; W) \approx d_i$ \emph{nos pontos amostrados}
    \item E \textbf{torcemos} para acertar onde não há amostra —
          a questão da \emph{generalização}, tema recorrente do curso
  \end{itemize}

  \begin{ideiabox}
    ``Treinar'' $=$ ``interpolar $N$ pares entrada--saída com uma função
    paramétrica flexível, esperando que ela se comporte bem fora deles''.
    A frase descreve do perceptron ao GPT.
  \end{ideiabox}
\end{frame}

% ------------------------------------------------------------
\begin{frame}{Para discutir em aula}
  \begin{quizbox}
    Como redes neurais são aproximadores universais, \emph{qualquer} rede, de
    \emph{qualquer} arquitetura, aproxima qualquer função com precisão
    arbitrária — e, com poucas amostras de treinamento, seus parâmetros podem
    ser perfeitamente determinados.

    \smallskip
    \opcoes{Verdadeiro \;/\; Falso}
  \end{quizbox}

  \pause

  \begin{respbox}
    \textbf{Falso, duas vezes.} (1) A aproximação universal exige
    \alert{capacidade mínima}: redes com poucos parâmetros podem ser
    aproximadores muito ruins. (2) Amostras finitas \emph{subdeterminam} a
    função: infinitas redes passam pelos mesmos pontos. O que fazemos é
    \emph{estimar} os parâmetros que melhor ajustam os dados — nunca
    ``recuperar perfeitamente'' a função-alvo desconhecida.
  \end{respbox}
\end{frame}

% ============================================================
\section{O perceptron de Rosenblatt (perspectiva histórica)}

% ------------------------------------------------------------
\begin{frame}{O perceptron original: unidade de limiar}
  \begin{columns}[T]
    \begin{column}{0.5\textwidth}
      \begin{itemize}\setlength{\itemsep}{0.4ex}
        \item Rosenblatt (1958): soma ponderada seguida de \textbf{degrau} —
              \smallskip
              \begin{center}
                $\hat{y} = 1$ se $\vw^{\!\top}\vx + b > 0$;\;
                $\hat{y} = 0$ caso contrário
              \end{center}
              \smallskip
        \item Truque de notação: o viés vira o peso $w_{N+1}$ de uma entrada
              fixa $x_{N+1} = 1$
        \item Com isso a fronteira $\vw^{\!\top}\vx = 0$ é um hiperplano
              \emph{pela origem} no espaço estendido
      \end{itemize}
    \end{column}
    \begin{column}{0.46\textwidth}
      \begin{center}
        \scalebox{0.60}{\perceptron}
      \end{center}
    \end{column}
  \end{columns}

  \begin{defbox}{Geometria do perceptron}
    $\vw$ é \textbf{ortogonal} ao hiperplano de decisão e aponta para a classe
    positiva. Aprender o perceptron $=$ encontrar esse vetor normal.
  \end{defbox}
\end{frame}

% ------------------------------------------------------------
\begin{frame}[fragile]{O algoritmo do perceptron}
  \begin{columns}[T]
    \begin{column}{0.47\textwidth}
      \begin{algbox}{Regra de Rosenblatt (\emph{online})}
        Com rótulos $d_i \in \{-1, +1\}$:
        \begin{itemize}\setlength{\itemsep}{0.2ex}
          \item Inicialize $\vw = \vet{0}$
          \item Percorra as amostras; se $d_i\,(\vw^{\!\top}\vx_i) \le 0$
                (erro), atualize
                \[ \vw \leftarrow \vw + d_i\, \vx_i \]
          \item Repita até não haver erros
        \end{itemize}
      \end{algbox}
      {\small Errou um positivo? Some $\vx_i$ (aproxima $\vw$ dele).
      Errou um negativo? Subtraia. Só isso.}
    \end{column}
    \begin{column}{0.51\textwidth}
\begin{lstlisting}[basicstyle=\ttfamily\scriptsize]
import numpy as np

def perceptron(X, y, epocas=100):
    # y em {-1, +1}; coluna do vies
    X = np.c_[X, np.ones(len(X))]
    w = np.zeros(X.shape[1])
    for _ in range(epocas):
        erros = 0
        for xi, yi in zip(X, y):
            if yi * (w @ xi) <= 0:
                w += yi * xi  # Rosenblatt
                erros += 1
        if erros == 0:
            break
    return w
\end{lstlisting}
    \end{column}
  \end{columns}
\end{frame}

% ------------------------------------------------------------
\begin{frame}{Quando funciona — e quando trava}
  \begin{teobox}{Convergência do perceptron (Novikoff)}
    Se as classes são \textbf{linearmente separáveis}, o algoritmo converge em
    tempo finito: partindo de $\vw = \vet{0}$, faz no máximo
    $\left(R/\gamma\right)^{2}$ atualizações, onde
    $R = \max_i \norm{\vx_i}$ e $\gamma$ é a \emph{margem} — a menor distância
    de uma amostra ao hiperplano separador ótimo.
  \end{teobox}

  \begin{itemize}\setlength{\itemsep}{0.3ex}
    \item A mesma $\gamma$ das SVMs: margens largas $\Rightarrow$ problema fácil
    \item Intuição do limitante: cada erro dá um passo de tamanho
          $\sim\!R$, e são precisos muitos passos de tamanho $\gamma$
          para desfazê-lo
  \end{itemize}

  \begin{alertabox}
    Se as classes \alert{não} são linearmente separáveis — o caso típico de
    dados reais, com ruído e sobreposição — não existe hiperplano perfeito e o
    algoritmo \textbf{nunca converge}: fica oscilando para sempre. Já aqui o
    modelo de 1958 esbarra na realidade.
  \end{alertabox}
\end{frame}

% ------------------------------------------------------------
\begin{frame}{Por que a regra do perceptron não treina um MLP}
  \begin{columns}[T]
    \begin{column}{0.56\textwidth}
      \begin{itemize}\setlength{\itemsep}{0.4ex}
        \item Fronteiras não lineares exigem \textbf{camadas ocultas}:
              cada neurônio oculto aprende \emph{uma} reta
        \item O dado só especifica \alert{entrada e saída da rede} — ninguém
              diz o que cada neurônio oculto deve responder
        \item A regra exigiria \emph{inventar rótulos intermediários}
              separáveis que, combinados, deem a fronteira certa —
              para \emph{todos} os neurônios ao mesmo tempo
      \end{itemize}
    \end{column}
    \begin{column}{0.40\textwidth}
      \begin{center}
      \scalebox{0.58}{%
      \begin{tikzpicture}[node distance=0.5cm and 1.1cm]
        \node[neuronio entrada] (x1) {$x_1$};
        \node[neuronio entrada, below=of x1] (x2) {$x_2$};
        \node[neuronio destacado, right=of x1] (h1) {\textbf{?}};
        \node[neuronio destacado, right=of x2] (h2) {\textbf{?}};
        \node[neuronio saida] (o) at ($(h1)!0.5!(h2)+(1.7,0)$) {$\hat{y}$};
        \draw[sinapse] (x1) -- (h1);
        \draw[sinapse] (x1) -- (h2);
        \draw[sinapse] (x2) -- (h1);
        \draw[sinapse] (x2) -- (h2);
        \draw[sinapse] (h1) -- (o);
        \draw[sinapse] (h2) -- (o);
        \node[below=0.28cm of h2, text width=3.2cm, align=center,
              icCinza, font=\small]
              {saídas desejadas \textbf{desconhecidas}};
      \end{tikzpicture}}
      \end{center}
    \end{column}
  \end{columns}

  \begin{alertabox}
    Rotular os neurônios internos é \textbf{busca combinatória} (exponencial).
    A constatação \alert{paralisou a área por mais de uma década}: o MLP
    parecia intreinável.
  \end{alertabox}
\end{frame}

% ------------------------------------------------------------
\begin{frame}{Curiosidade histórica: ADALINE, MADALINE e a regra delta}
  \framesubtitle{Widrow \& Hoff (1960) — o detalhe que sobreviveu foi a ideia de gradiente}

  \begin{itemize}\setlength{\itemsep}{0.3ex}
    \item \textbf{ADALINE:} mesmo perceptron, outra regra de aprendizado —
          minimiza o erro quadrático sobre a \emph{pré-ativação} $\campo$
          (contínua!), não sobre a saída binária
    \item \textbf{MADALINE:} versão multicamada com heurística gulosa (inverte
          a saída do neurônio ``mais em dúvida'') — não escalou; nota de rodapé
  \end{itemize}

  \[
    \erro = \tfrac{1}{2}\,(d - \campo)^2
    \qquad\Longrightarrow\qquad
    \mdestaque{icMagenta}{\vw \leftarrow \vw + \taxa\,(d - \campo)\,\vx}
  \]

  \begin{ideiabox}
    A \textbf{regra delta} (LMS, Widrow--Hoff) \emph{é} a descida de gradiente
    do erro quadrático em um neurônio linear. Generalizada para qualquer
    ativação diferenciável e encadeada camada a camada, ela vira exatamente a
    \alert{retropropagação} — o algoritmo que treina toda rede moderna. De
    1960 só descartamos o degrau; o gradiente ficou.
  \end{ideiabox}
\end{frame}

% ------------------------------------------------------------
\begin{frame}{Para discutir em aula}
  \begin{quizbox}
    Em um MLP cuja fronteira-alvo é um polígono, todos os neurônios já foram
    corretamente ajustados, exceto \emph{um} neurônio oculto. Como os demais
    estão certos, podemos usar a regra do perceptron para aprender esse último
    neurônio diretamente dos dados de treinamento.

    \smallskip
    \opcoes{Verdadeiro \;/\; Falso}
  \end{quizbox}

  \pause

  \begin{respbox}
    \textbf{Falso.} Os rótulos do conjunto de treinamento descrevem a saída da
    \emph{rede}, não a saída desejada \emph{daquele neurônio} — e, vistos por
    ele, os dados nem sequer são linearmente separáveis. Seria preciso
    \emph{rerrotular} instâncias para criar o subproblema linear correto, e o
    dado não diz quais. É a versão em miniatura do problema combinatório que
    inviabiliza a regra do perceptron em redes profundas.
  \end{respbox}
\end{frame}

% ============================================================
\section{Diferenciabilidade: a virada moderna}

% ------------------------------------------------------------
\begin{frame}{O verdadeiro vilão: derivada nula}
  \begin{columns}[T]
    \begin{column}{0.55\textwidth}
      \begin{itemize}\setlength{\itemsep}{0.4ex}
        \item O erro de classificação é \textbf{binário}: deslocar levemente a
              fronteira quase nunca muda a contagem
        \item O degrau é \textbf{plano} (derivada $0$) e não diferenciável
              no salto
        \item Mexer um pouquinho em um peso \alert{não altera o erro} — não há
              sinal de \emph{qual direção} melhora; o erro da rede vira uma
              função em platôs: sem gradiente, sem bússola
      \end{itemize}
    \end{column}
    \begin{column}{0.41\textwidth}
      \begin{center}
      \scalebox{0.85}{%
      \begin{tikzpicture}
        \draw[->, icCinza] (-2.1,0) -- (2.3,0) node[below] {\small $\campo$};
        \draw[->, icCinza] (0,-0.25) -- (0,1.45);
        \draw[very thick, icCarmim] (-2.1,0) -- (0,0);
        \draw[very thick, icCarmim] (0,1) -- (2.1,1);
        \fill[icCarmim] (0,1) circle (1.6pt);
        \draw[icCarmim, fill=white] (0,0) circle (1.6pt);
        \node[icCarmim, font=\small] at (1.4,1.25) {degrau};
        \node[icCinza, font=\scriptsize, align=center] at (0.1,-0.65)
          {derivada $0$ em toda parte;\\ indefinida no salto};
      \end{tikzpicture}}
      \end{center}
    \end{column}
  \end{columns}

  \begin{alertabox}
    O obstáculo nunca foi ``multicamadas''; foi \alert{a ausência de gradiente}.
    Com derivadas, o problema vira otimização contínua — tratável mesmo com
    bilhões de parâmetros.
  \end{alertabox}
\end{frame}

% ------------------------------------------------------------
\begin{frame}{Os dois requisitos da aprendizagem moderna}
  \begin{defbox}{Requisitos de treinabilidade}
    \textbf{(1) Ativações diferenciáveis} — a saída de cada neurônio varia
    suavemente com pesos e entradas; \textbf{(2) perda diferenciável} — o
    desajuste rede\,$\times$\,alvo varia suavemente com a saída da rede.
  \end{defbox}

  \begin{itemize}\setlength{\itemsep}{0.3ex}
    \item Degrau $\to$ funções suaves: mudanças de peso produzem mudanças
          \emph{mensuráveis} na saída
    \item ``Conte os erros'' $\to$ \alert{divergência contínua}, um
          \emph{proxy} do erro de classificação
    \item Pela \textbf{regra da cadeia}, a rede \emph{inteira} vira
          diferenciável em cada parâmetro — o que a retropropagação
          explora (Aula 5)
  \end{itemize}

  \begin{ideiabox}
    O coração do curso: \emph{deep learning é a arte de manter tudo
    diferenciável}, para que o gradiente flua da perda a cada parâmetro.
  \end{ideiabox}
\end{frame}

% ------------------------------------------------------------
\begin{frame}{Ativações modernas}
  \begin{columns}[T]
    \begin{column}{0.31\textwidth}
      \begin{center}
      \scalebox{0.72}{%
      \begin{tikzpicture}
        \draw[->, icCinza] (-2.2,0) -- (2.4,0) node[below] {\small $\campo$};
        \draw[->, icCinza] (0,-0.25) -- (0,1.5);
        \draw[thick, icMagenta]
          plot[domain=-2.2:2.2, samples=60] (\x, {1/(1+exp(-2.4*\x))});
        \node[icMagenta, font=\small] at (1.2,1.35) {$\sig(\campo)$};
      \end{tikzpicture}}
      \end{center}
      {\small\textbf{Sigmoide} $\sig(\campo)=\dfrac{1}{1+e^{-\campo}}$\\
      saída em $(0,1)$; interpretação probabilística}
    \end{column}
    \begin{column}{0.31\textwidth}
      \begin{center}
      \scalebox{0.72}{%
      \begin{tikzpicture}
        \draw[->, icCinza] (-2.2,0) -- (2.4,0) node[below] {\small $\campo$};
        \draw[->, icCinza] (0,-0.25) -- (0,1.5);
        \draw[thick, icAzulClaro] (-2.2,0) -- (0,0);
        \draw[thick, icAzulClaro] (0,0) -- (1.45,1.45);
        \node[icAzulClaro, font=\small] at (0.8,1.35) {$\relu(\campo)$};
      \end{tikzpicture}}
      \end{center}
      {\small\textbf{ReLU} $\max(0,\campo)$\\
      padrão em redes profundas; barata, gradiente $1$ no lado ativo}
    \end{column}
    \begin{column}{0.31\textwidth}
      \begin{center}
      \scalebox{0.72}{%
      \begin{tikzpicture}
        \draw[->, icCinza] (-2.2,0) -- (2.4,0) node[below] {\small $\campo$};
        \draw[->, icCinza] (0,-1.3) -- (0,1.5);
        \draw[thick, icVerde]
          plot[domain=-2.2:2.2, samples=60] (\x, {tanh(1.2*\x)});
        \node[icVerde, font=\small] at (1.35,1.3) {$\tanh(\campo)$};
      \end{tikzpicture}}
      \end{center}
      {\small\textbf{Tanh} — sigmoide centrada em $0$; e primas modernas:
      \emph{softplus}, GELU, SiLU\ldots}
    \end{column}
  \end{columns}

  \medskip
  \begin{itemize}\setlength{\itemsep}{0.2ex}
    \item Diferenciáveis \emph{em quase todo lugar} basta: no ``bico'' da ReLU
          usamos uma \textbf{subderivada} (qualquer valor em $[0,1]$) — e o
          treinamento nem percebe
  \end{itemize}
\end{frame}

% ------------------------------------------------------------
\begin{frame}{A sigmoide é especial: nasce a regressão logística}
  \begin{columns}[T]
    \begin{column}{0.55\textwidth}
      \begin{defbox}{Regressão logística}
        Um perceptron com ativação sigmoide computa
        \[
          P(y = 1 \mid \vx) \;=\; \sig\!\left(\vw^{\!\top}\vx + b\right),
        \]
        a \textbf{probabilidade} da classe $1$; decide-se $\hat{y}=1$
        se ela excede $0{,}5$.
      \end{defbox}
      \begin{itemize}\setlength{\itemsep}{0.3ex}
        \item Em dados reais as classes se \emph{sobrepõem}: a fração de
              positivos cresce suavemente de $0$ a $1$ — o formato da sigmoide
        \item A saída vira \alert{grau de confiança}, não um veredito
      \end{itemize}
    \end{column}
    \begin{column}{0.41\textwidth}
      \begin{center}
      \scalebox{0.70}{%
      \begin{tikzpicture}
        \draw[->, icCinza] (-2.3,0) -- (2.5,0) node[below] {\small $x$};
        \draw[->, icCinza] (-2.3,0) -- (-2.3,1.6) node[left] {\small $y$};
        \foreach \p in {-2.1,-1.8,-1.55,-1.1,-0.7,-0.2,0.3}
          \fill[icAzulClaro] (\p,0) circle (1.5pt);
        \foreach \p in {-0.5,0.05,0.5,0.9,1.3,1.7,2.1,2.3}
          \fill[icCarmim] (\p,1.3) circle (1.5pt);
        \draw[thick, icMagenta]
          plot[domain=-2.3:2.4, samples=60]
          (\x, {1.3/(1+exp(-2.6*(\x-0.1)))});
        \node[icCinza, font=\scriptsize, align=center] at (0.2,-0.55)
          {classes sobrepostas: nenhum limiar\\ é perfeito, mas $P(y{=}1)$
           cresce suave};
      \end{tikzpicture}}
      \end{center}
    \end{column}
  \end{columns}

\end{frame}

% ------------------------------------------------------------
\begin{frame}[fragile]{Diferenciabilidade em código: \texttt{autograd}}

\begin{lstlisting}
import torch

x = torch.tensor([1.0, 2.0])
w = torch.tensor([0.5, -0.3], requires_grad=True)   # parametros:
b = torch.tensor(0.1, requires_grad=True)           # pediremos gradiente

z = w @ x + b               # campo local induzido (diferenciavel)
y = torch.sigmoid(z)        # ativacao suave (diferenciavel)
d = torch.tensor(1.0)       # alvo
loss = (y - d)**2           # divergencia (diferenciavel)
loss.backward()             # regra da cadeia, automatica
print(w.grad, b.grad)       # dloss/dw e dloss/db
\end{lstlisting}

  \begin{ideiabox}
    Troque \texttt{sigmoid} por um degrau e \texttt{w.grad} vira zero em toda
    parte: o otimizador fica cego. Ativações suaves são \emph{a} condição
    para o \texttt{backward()} existir.
  \end{ideiabox}
\end{frame}

% ------------------------------------------------------------
\begin{frame}{Para discutir em aula}
  \begin{quizbox}
    Com ativação degrau, a derivada do erro de classificação em relação ao
    limiar indica em que direção movê-lo. Já com ativação sigmoide, deslocar a
    função altera a distância total entre saídas e alvos mesmo quando o erro
    de classificação não muda.

    \smallskip
    \opcoes{Só a 1\textsuperscript{a} é verdadeira \;/\;
            Só a 2\textsuperscript{a} é verdadeira \;/\;
            Ambas \;/\; Nenhuma}
  \end{quizbox}

  \pause

  \begin{respbox}
    \textbf{Só a 2\textsuperscript{a}.} Com degrau, o erro é constante por
    trechos: a derivada é $0$ (ou inexistente) e \emph{não} aponta direção
    alguma. Com sigmoide, a distância contínua aos alvos muda a cada
    deslocamento — e o seu gradiente diz para onde ir, \emph{mesmo entre}
    duas configurações com a mesma contagem de erros. É exatamente por isso
    que otimizamos um \alert{proxy contínuo}, e não o erro $0/1$.
  \end{respbox}
\end{frame}

% ============================================================
\section{Minimização de risco empírico}

% ------------------------------------------------------------
\begin{frame}{Divergência: medindo o desajuste de forma diferenciável}
  \begin{defbox}{Função de divergência}
    $\dive(f(\vx; W),\, g(\vx))$ quantifica o desajuste entre a saída da rede
    e o alvo, satisfazendo:
    \begin{itemize}\setlength{\itemsep}{0.15ex}
      \item $\dive(f, g) = 0$ quando $f(\vx; W) = g(\vx)$;
      \item $\dive(f, g) > 0$ quando diferem;
      \item $\dive$ é \textbf{diferenciável em relação a $f$} (e portanto,
            pela cadeia, em relação a $W$).
    \end{itemize}
  \end{defbox}

  \begin{itemize}\setlength{\itemsep}{0.3ex}
    \item Para classificação, a divergência \alert{não é} o erro $0/1$ —
          é um substituto suave dele (veremos: entropia cruzada)
    \item Para regressão, o clássico: $\dive = \norm{f(\vx;W) - g(\vx)}^2$
    \item A escolha da divergência é decisão de modelagem tão importante
          quanto a arquitetura
  \end{itemize}
\end{frame}

% ------------------------------------------------------------
\begin{frame}{Risco esperado, risco empírico}
  O que gostaríamos de minimizar — o \textbf{risco esperado} sobre todo o
  domínio de entrada:
  \[
    R(W) \;=\; \E_{\vx}\!\left[\, \dive\!\big(f(\vx; W),\, g(\vx)\big) \,\right]
  \]

  O que \emph{conseguimos} calcular — a média sobre as $N$ amostras:
  \[
    \mdestaque{icMagenta}{\hat{R}(W) \;=\;
      \frac{1}{N} \sum_{i=1}^{N}
      \termo{\dive\!\big(f(\vx_i; W),\, d_i\big)}{desajuste na amostra $i$}}
  \]

  \begin{defbox}{Minimização de risco empírico (ERM)}
    Treinar a rede é resolver
    \(\;\hat{W} = \argmin_{W}\, \hat{R}(W)\;\)
    — minimizar a \emph{estimativa empírica} do risco, na esperança de que ela
    aproxime bem o risco verdadeiro.
  \end{defbox}

  {\small\color{icCinza} Na prática chamamos $\hat{R}$ de \emph{loss}
  (perda). Para um conjunto de treinamento fixo, a perda é função apenas
  de $W$.}
\end{frame}

% ------------------------------------------------------------
\begin{frame}{As divergências do dia a dia moderno}
  \begin{columns}[T]
    \begin{column}{0.48\textwidth}
      \begin{tcolorbox}[iccaixa=icCiano, title={Regressão: erro quadrático}]
        \[
          \dive = \norm{\hat{\vy}_i - \vy_i}^2
        \]
        {\small \texttt{nn.MSELoss} — suave, penaliza quadraticamente;
        ligado à verossimilhança gaussiana.}
      \end{tcolorbox}
    \end{column}
    \begin{column}{0.48\textwidth}
      \begin{tcolorbox}[iccaixa=icMagenta, title={Classificação: entropia cruzada}]
        \[
          \dive = -\!\sum_c d_{i,c} \log \hat{y}_{i,c}
        \]
        {\small \texttt{nn.CrossEntropyLoss} / \texttt{nn.BCEWithLogitsLoss}
        — proxy suave do erro $0/1$; deriva da leitura probabilística
        da sigmoide/softmax.}
      \end{tcolorbox}
    \end{column}
  \end{columns}

  \medskip
  \begin{itemize}\setlength{\itemsep}{0.3ex}
    \item As duas são diferenciáveis, convexas \emph{na saída da rede} e têm
          gradientes bem-comportados — por isso dominam a prática
    \item Deduziremos a entropia cruzada com calma quando tratarmos de saídas
          probabilísticas; hoje basta o papel: \alert{divergência no ERM}
  \end{itemize}
\end{frame}

% ------------------------------------------------------------
\begin{frame}[fragile]{ERM em PyTorch: o esqueleto de todo treinamento}
\begin{lstlisting}[basicstyle=\ttfamily\scriptsize]
import torch, torch.nn as nn

model = nn.Sequential(              # f( . ; W): MLP diferenciavel
    nn.Linear(2, 16), nn.ReLU(),
    nn.Linear(16, 1))
criterio = nn.BCEWithLogitsLoss()   # div(): proxy suave do erro 0/1
otim = torch.optim.SGD(model.parameters(), lr=0.1)
for epoca in range(1000):
    saida = model(X).squeeze()      # f(X_i; W) para todo i
    risco = criterio(saida, d)      # R^(W): media das divergencias
    otim.zero_grad()
    risco.backward()                # gradientes de R^ em cada peso
    otim.step()                     # ajusta W na direcao que reduz R^
\end{lstlisting}

  \begin{ideiabox}
    Mapeie o laço na teoria: \texttt{model} $=f(\cdot;W)$, \texttt{criterio}
    $=\dive$, \texttt{risco} $=\hat{R}(W)$, \texttt{backward} $=$ regra da
    cadeia, \texttt{step} $=$ otimização. Do MLP aos LLMs, é este o programa.
  \end{ideiabox}
\end{frame}

% ------------------------------------------------------------
\begin{frame}{Treinar $=$ otimizar: a ponte para a próxima aula}
  \begin{itemize}\setlength{\itemsep}{0.4ex}
    \item Chegamos a um enunciado puramente matemático: dado $\D$, minimizar
          $\hat{R}(W)$ em relação a $W$ — uma instância de
          \textbf{minimização de funções}
    \item Como tudo é diferenciável, o gradiente
          $\deriv{\hat{R}}{W}$ existe e aponta a direção de
          \emph{maior crescimento} do risco — andaremos no sentido oposto
  \end{itemize}

  \[
    \mdestaque{icMagenta}{W \;\leftarrow\; W \;-\;
      \termo{\taxa}{taxa de aprendizado} \cdot
      \termo{\deriv{\hat{R}}{W}}{gradiente do risco empírico}}
  \]

  \begin{itemize}\setlength{\itemsep}{0.4ex}
    \item \textbf{Aula 4:} descida de gradiente em detalhe — o que é o
          gradiente em muitas dimensões, mínimos locais, taxa de aprendizado
    \item \textbf{Aula 5:} como calcular esse gradiente \emph{eficientemente}
          em uma rede profunda — a retropropagação
  \end{itemize}
\end{frame}

% ------------------------------------------------------------
\begin{frame}{Exercício}
  \begin{exbox}{1 — Da regra delta ao gradiente}
    Considere um neurônio linear $\hat{y} = \vw^{\!\top}\vx$ com perda
    $\erro = \tfrac{1}{2}(d - \hat{y})^2$ para uma única amostra
    $(\vx, d)$.
    \begin{itemize}\setlength{\itemsep}{0.2ex}
      \item[(a)] Calcule $\deriv{\erro}{\vw}$ e mostre que a atualização por
                 gradiente, com passo $\taxa$, é exatamente a regra delta
                 de Widrow--Hoff vista hoje.
      \item[(b)] Repita com $\hat{y} = \sig(\vw^{\!\top}\vx)$, usando
                 $\sig'(\campo) = \sig(\campo)\,(1 - \sig(\campo))$. Que fator novo
                 aparece na atualização?
      \item[(c)] Verifique (b) numericamente com o trecho de
                 \texttt{autograd} do slide anterior, comparando
                 \texttt{w.grad} com sua fórmula.
    \end{itemize}
  \end{exbox}
\end{frame}

% ------------------------------------------------------------
\begin{frame}{Resumo da aula}
  \begin{itemize}\setlength{\itemsep}{0.45ex}
    \item Aprender a rede $=$ estimar os parâmetros $W$ de $f(\vx; W)$ a partir
          de \textbf{amostras} $(\vx_i, d_i)$ da função-alvo
    \item A regra do perceptron aprende \emph{um} separador linear em tempo
          finito (classes separáveis) — mas em redes multicamadas exige
          rótulos intermediários desconhecidos: \alert{busca combinatória};
          ADALINE/MADALINE ficaram como curiosidade histórica
    \item O desbloqueio moderno foi tornar tudo \textbf{diferenciável}:
          ativações suaves (sigmoide, ReLU, \ldots) $+$ divergência contínua
          como proxy do erro — o \texttt{autograd} é a materialização disso
    \item \textbf{ERM:} minimizar o risco empírico
          $\hat{R}(W) = \frac{1}{N}\sum_i \dive(f(\vx_i; W), d_i)$ — o laço de
          treinamento do PyTorch é ERM ao pé da letra
    \item Treinar virou \textbf{otimização de função diferenciável};
          próxima aula: descida de gradiente
  \end{itemize}
\end{frame}

\end{document}
